\section{Conclusion \& Recommendations}

\subsection{Project Outcomes}

The project successfully delivered a fully functional Hospital Management System that integrates administrative efficiency with secure data persistence. Key achievements include:
\begin{itemize}
    \item \textbf{Three-Tier Architecture:} Established a clear separation between the web frontend, Python/Flask backend, and MySQL database layer for improved maintainability.
    \item \textbf{RESTful API:} Developed a robust backend with five resource endpoints handling CRUD operations through parameterized queries to ensure security against SQL injection.
    \item \textbf{Database Automation:} Implemented Stored Procedures, Triggers, Views, and Indexes to enforce business rules at the database level.
    \item \textbf{Comprehensive Reporting:} Delivered a reporting module that provides automated insights into daily schedules, financial summaries, and departmental workloads with minimal latency.
\end{itemize}

\subsection{Lessons Learned}

The development process provided critical insights into the design of data-driven applications:
\begin{itemize}
    \item \textbf{Logic Delegation:} While moving business logic to stored procedures and triggers ensures high data integrity, it increases debugging complexity as errors must be propagated carefully from the database to the application layer.
    \item \textbf{Connection Management:} The per-request connection pattern used in \texttt{database.py} is reliable for small-scale deployment but introduces overhead that limits scalability under high concurrency.
    \item \textbf{Data Serialization:} Systematic type normalization for Date and Decimal types is essential for ensuring consistent JSON responses across all API endpoints.
\end{itemize}

\subsection{Future Work}

To enhance the system for enterprise-level use, the following improvements are proposed:
\begin{itemize}
    \item \textbf{Technical Hardening:} Implementing connection pooling (SQLAlchemy) for high concurrency, moving credentials to environment variables, and migrating to a cloud environment for remote accessibility.
    \item \textbf{Frontend Evolution:} Shifting to a component-based framework like React to handle more complex UI state management.
    \item \textbf{Functional Expansion:} Developing modules for Pharmacy and Medication management, Electronic Medical Records (EMR), and third-party insurance integration to support end-to-end hospital workflows.
\end{itemize}